\begin{thebibliography}{99}
\bibitem{klimesch1999} Klimesch W. EEG alpha and theta oscillations reflect cognitive and memory performance: a review and analysis. Brain Research Reviews. 1999;29(2-3):169-195.
\bibitem{gevins1997} Gevins A, Smith ME, McEvoy L, Yu D. High-resolution EEG mapping of cortical activation related to working memory: effects of task difficulty, type of processing, and practice. Cerebral Cortex. 1997;7(4):374-385.
\bibitem{sridhar2022} Sridhar S, Manian V, Rao RPN. Theta and alpha ratios as objective correlates of cognitive workload: a cross-task replication study. Journal of Neural Engineering. 2022;19(4):046033.
\bibitem{zyma2019} Zyma I, Tukaev S, Seleznov I, Kiyono K, Popov A, Chernykh M, Shpenkov O. Electroencephalograms during Mental Arithmetic Task Performance. Data. 2019;4(1):14.
\bibitem{lotte2018} Lotte F, Bougrain L, Cichocki A, Clerc M, Congedo M, Rakotomamonjy A, Yger F. A review of classification algorithms for EEG-based brain-computer interfaces: a 10 year update. Journal of Neural Engineering. 2018;15(3):031005.
\bibitem{physionet_eegmat} Zyma I, Tukaev S, Seleznov I, Kiyono K, Popov A, Chernykh M, Shpenkov O. EEG During Mental Arithmetic Tasks [Dataset]. PhysioNet. 2019. https://doi.org/10.13026/C2JQ1P.
\bibitem{stew_dataset} Lim WL, Sourina O, Wang LP. STEW: Simultaneous task EEG workload data set. IEEE Transactions on Neural Systems and Rehabilitation Engineering. 2018;26(11):2106-2114.
\bibitem{openneuro_ds007262} Barras M, Booth L. Cognitive Workload 8-level Arithmetic Task. OpenNeuro DS007262. 2026. https://doi.org/10.18112/openneuro.ds007262.v1.0.6.
\bibitem{mne} Gramfort A, Luessi M, Larson E, Engemann DA, Strohmeier D, Brodbeck C, Goj R, Jas M, Brooks T, Parkkonen L, Hamalainen M. MEG and EEG data analysis with MNE-Python. Frontiers in Neuroscience. 2013;7:267.
\bibitem{jayaram2016} Jayaram V, Alamgir M, Altun Y, Scholkopf B, Grosse-Wentrup M. Transfer learning in brain-computer interfaces. IEEE Computational Intelligence Magazine. 2016;11(1):20-31.
\bibitem{hakimi2023} Hakimi N, Jodeiri A, Setarehdan SK. Mental arithmetic task classification using convolutional neural networks on EEG signals. Biomedical Signal Processing and Control. 2023;79:104156.
\bibitem{roy2016} Roy RN, Charbonnier S, Campagne A, Bonnet S. Efficient mental workload estimation using task-independent EEG features. Journal of Neural Engineering. 2016;13(2):026019.
\bibitem{jebelli2019} Jebelli H, Hwang S, Lee S. EEG-based workers' mental workload classification using deep learning. Automation in Construction. 2019;106:102876.
\bibitem{nguyen2025} Nguyen PK, Huynh QL. A machine learning approach in EEG-based assessment of cognitive load by mental arithmetic tasks. Journal of Physics: Conference Series. 2025;2949:012004.
\bibitem{borghini2014} Borghini G, Astolfi L, Vecchiato G, Mattia D, Babiloni F. Measuring neurophysiological signals in aircraft pilots and car drivers for the assessment of mental workload, fatigue and drowsiness. Neuroscience \& Biobehavioral Reviews. 2014;44:58-75.
\bibitem{cavanagh2014} Cavanagh JF, Shackman AJ. Frontal midline theta reflects anxiety and cognitive control: a meta-analytic review. Journal of Physiology-Paris. 2014;108(4-6):251-258.
\bibitem{sweller1988} Sweller J. Cognitive load during problem solving: effects on learning. Cognitive Science. 1988;12(2):257-285.
\bibitem{hart1988} Hart SG, Staveland LE. Development of NASA-TLX (Task Load Index): results of empirical and theoretical research. Advances in Psychology. 1988;52:139-183.
\bibitem{gramfort2013} Gramfort A, Luessi M, Larson E, Engemann DA, Strohmeier D, Brodbeck C, Goj R, Jas M, Brooks T, Parkkonen L, Hamalainen M. MEG and EEG data analysis with MNE-Python. Frontiers in Neuroscience. 2013;7:267.
\bibitem{sklearn} Pedregosa F, Varoquaux G, Gramfort A, Michel V, Thirion B, Grisel O, et al. Scikit-learn: Machine learning in Python. Journal of Machine Learning Research. 2011;12:2825-2830.
\bibitem{zander2011} Zander TO, Kothe C. Towards passive brain-computer interfaces: applying brain-computer interface technology to human-machine systems in general. Journal of Neural Engineering. 2011;8(2):025005.
\bibitem{vidaurre2011} Vidaurre C, Sannelli C, Muller KR, Blankertz B. Co-adaptive calibration to improve BCI performance. NeuroImage. 2011;55(4):1443-1456.
\bibitem{congedo2017} Congedo M, Barachant A, Bhatia A. Riemannian geometry for EEG-based brain-computer interfaces: a primer and a review. Brain-Computer Interfaces. 2017;4(3):155-174.
\bibitem{baldwin2013} Baldwin CL, Penaranda BN. Adaptive training using an artificial neural network and EEG metrics for within- and cross-task workload classification. NeuroImage. 2013;59(1):48-56.
\bibitem{goldberger2000} Goldberger AL, Amaral LAN, Glass L, Hausdorff JM, Ivanov PC, Mark RG, et al. PhysioBank, PhysioToolkit, and PhysioNet: Components of a new research resource for complex physiologic signals. Circulation. 2000;101(23):e215-e220.
\bibitem{wong2011} Wong B. Points of view: Color blindness. Nature Methods. 2011;8(6):441.
\bibitem{zanini2018} Zanini P, Congedo M, Jutten C, Said S, Berthoumieu Y. Transfer learning: a Riemannian geometry framework with applications to brain-computer interfaces. IEEE Transactions on Biomedical Engineering. 2018;65(5):1107-1116.
\end{thebibliography}
